\documentclass[a4paper,twoside]{article}
\usepackage[T1]{fontenc}
\usepackage[bahasa]{babel}
\usepackage{graphicx}
\usepackage{graphics}
\usepackage{float}
\usepackage[cm]{fullpage}
\pagestyle{myheadings}
\usepackage{etoolbox}
\usepackage{setspace} 
\usepackage{lipsum} 
\usepackage{amsmath}
\usepackage{booktabs}
\usepackage{forest}
\usepackage{longdivision}
\setlength{\headsep}{30pt}
\usepackage[inner=2cm,outer=2.5cm,top=2.5cm,bottom=2cm]{geometry} %margin
% \pagestyle{empty}

\makeatletter
\renewcommand{\@maketitle} {\begin{center} {\LARGE \textbf{ \textsc{\@title}} \par} \bigskip {\large \textbf{\textsc{\@author}} }\end{center} }
\renewcommand{\thispagestyle}[1]{}

\onehalfspacing

\begin{document}
\section{Stat-DSM}
Tujuan utama dari Stat-DSM adalah jika diberikan dua kelompok lintasan, Stat-DSM ingin untuk menemukan pola pergerakan dalam bentuk sub-lintasan yang lebih mirip dengan sub-lintasan di satu kelompok, namun kurang mirip dengan sub-lintasan di kelompok lainnya.
Stat-DSM ingin memisahkan dua kelas lintasan dimana setiap sub-lintasan pada setiap kelas memiliki perbedaan yang cukup sehingga dapat dinyatakan bahwa perbedaan tersebut merupakan signifikan secara statistik dan tidak terjadi hanyak karena acak.

\section{Lintasan}
Sebuah lintasan diisi dengan kumpulan titik yang bersifat kontinu $T_i = (p_i^1,p_i^2,...,p_i^n)$.
Dengan setiap titik dinotasikan dengan $p_i^t = (x_i^t,y_i^t)$ yang merupakan koordinat x dan y dari setiap titik pada waktu.
Setiap lintasan memiliki kelas atau label $Y_i\in\{0,1\}$ dimana $Y_i = 1$ menandakan sukses dan $Y_i = 0$ menandakan tidak sukses.

\section{Sub-Lintasan}
Sub-lintasan adalah bagian kontinu dari suatu lintasan.
Anggap kita memiliki sebuah lintasan \[T_i = (p_i^1,p_i^2,p_i^3,p_i^4,p_i^5,p_i^6,p_i^7,p_i^8)\]sebuah sub-lintasan dapat berupa $S_i = (p_i^3,p_i^4,p_i^5,p_i^6,p_i^7)$.

\section{Distance Function}
Pada Stat-DSM, distance function dipergunakan untuk mencari kemiripan antar dua sub-lintasan. 
Algoritma melakukan ini dengan cara membandingkan jarak koordinat antar dua sub-lintasan.
Algoritma pertama mencari sub-lintasan yang sama dengan menggunakan euclidean distance untuk mencari jarak antar dua titik.
Anggap kita memiliki dua titik yaitu $(x_1, y_1)$ dan $(x_2, y_2)$, maka \[ d = \sqrt{(x_2 - x_1)^2 + (y_2 - y_1)^2} \]
Lalu setelah setiap titik dicari euclidean ditance masing-masing, algoritma akan mengurut hasilnya dari yang peling besar hingga yang paling kecil.
Setelah diurut, kita akan menggunakan paramater K untuk menentukan menggunakan berapa distance terbesar untuk kalkulasi selanjutnya yaitu menggunakan \[ \operatorname{dist}_K(A,B)=\frac{1}{K}\sum_{k=1}^{K}d^{(k)} \] yang berarti jika $K = 1$, maka ambil euclidean distance yang terbesar, dan jika diatas 1 maka contohnya adalah $K = 2$ maka ambil 2 euclidean distance terbesar, tambahkan keduanya dan bagi dengan $K$

\section{Contingency Table}
Sebuah tabel dalam format matriks yang menampilkan distribusi frekuensi dari dua variabel kategorikal atau lebih.
Anggap kita menemukan untuk label kelas positif 80 mengandung $S$ dan 20 tidak dan untuk label kelas negatif 30 mengandung $S$ dan 70 tidak.
Sehingga didapatkannya 

\begin{table}[ht]
	\centering
	\caption{A 2$\times$2 Contingency Table}
	\label{tab:contingency_example}
	\begin{tabular}{rccc}
		\toprule
		& \multicolumn{2}{c}{\textbf{Mengandung}} & \\
		\cmidrule(lr){2-3}
		\textbf{Kelas} & $S$ & $\neg S$ & \textbf{Total} \\
		\midrule
		Positif & 80 & 20 & 100 \\
		Negatif & 30 & 70 & 100 \\
		\midrule
		\textbf{Total} & 110 & 90 & 200 \\
		\bottomrule
	\end{tabular}
\end{table} 

\section{Fisher's Exact Test}
Fisher's Exact Test merupakan sebuah metode untuk memeriksa apakah terdapat asosiasi (hubungan atau ketergantungan) yang signifikan antara dua variabel kategorikal.
Dengan menggunakan contingency table 

\begin{table}[ht]
	\centering
	\begin{tabular}{rccc}
		\textbf{Label} & $S$ & $\neg S$ \\
		\midrule
		Positif & a & b \\
		Negatif & c & d \\
	\end{tabular}
\end{table}

Dapat ditarik probabilitas mendapatkan data tersebut murni sebagai kebetulan dengan \[ p\text{-value} = \frac{\binom{a+b}{a}\binom{c+d}{c}}{\binom{a+b+c+d}{a+c}} \]

\section{Candidate Tree}
Pohon yang mengandung sub-lintasan yang akan diproses.
Setiap node pertama pada branch diisi dengan sub-lintasan ukuran $min_l$ dengan node dibawahnya merupakan sebuah perpanjangan sub-lintasan tersebut.
Penelusuran pohon dilakukan secara Depth-First-Search, yang berarti algoritma memproses sebuah sub-lintasan dengan ukuran $min_l$ dan mencoba untuk memperpanjang sub-lintasan tersebut sampai tidak bisa lagi atau $p$-value lebih besar dari threshold lalu lanjut ke sub-lintasan panjang $min_l$ lainnya.
Selama penelusuran pohon, dilakukan pula koreksi terhadap threshold dan threshold tersebut kemudian digunakan untuk menentukan branch mana dan di bagian mana pada pohon yang dapat di \textit{prune}
Contoh candidate tree :

\begin{figure}[h]
	\centering
	\begin{forest}
		[Root
			[$T_1: (A \rightarrow B)$
				[$T_1: (A \rightarrow B \rightarrow C)$
                    [$T_1: (A \rightarrow B \rightarrow C \rightarrow D)$]
                ]
			]
			[$T_1: (B \rightarrow C)$
				[$T_1: (B \rightarrow C \rightarrow D)$]
			]
			[$T_2: (A \rightarrow B)$
				[$T_2: (A \rightarrow B \rightarrow C)$]
			]
		]
	\end{forest}
	\caption{Contoh \textit{Candidate Tree}}
\end{figure}

\section{Westfall-Young Permutation Correction Method}
Sebuah metode statistik untuk menyesuaikan threshold yang digunakan ketika melakukan pengujian hipotesis pada ribuan hingga jutaan kandidat sub-lintasan secara bersamaan.
Cara kerjanya dalam algoritma Stat-DSM adalah pertama terdapat list $p$-value terkecil untuk setiap permutasi dari antara seluruh node. 
List tersebut berupa array dengan setiap index pada array tersebut dibedakan berdasarkan cara permutasi yang berbeda-beda untuk label dan isinya merupakan $p$-value terkecil untuk permutasi tersebut.
List array ini akan dinotasikan sebagai $p\_min(b)$ dimana b merupakan index dari array tersebut.
Kemudian pada setiap perhitungan node pada tree, akan dilakukan kalkulasi seluruh index b pada $p\_min$ yang pertama melihat apakah pada $p\_min(b)$ nilai $p$-value lebih besar dari nilai $p$-value dari node yang sekarang lagi diproses dan jika iya, maka lakukan perhitungan $p$-value menggunakan permutasi yang sama dan jika didapatkan $p$-value lebih kecil dari pada $p\_min(b)$ maka ubah isi $p\_min(b)$ menjadi nilai $p$-value dari permutasi tersebut.

\section{Input}
Input dari Stat-DSM meliputi 
\begin{enumerate}
    \item $T$ yang merupakan dataset lintasan 
    \item $g$ yang merupakan label untuk dataset lintasan
    \item $\varepsilon$ atau distance threshold untuk kalkulasi distance 
    \item $min_l$ untuk menentukan panjang minimum sebuah sub-lintasan
    \item $K$ atau Top-K-Max yang menentukan berapa nilai yang diambil untuk menghitung distance
    \item $B$ untuk membatas jumlah permutasi dari label untuk metode Westfall-Young. 
    \item $\alpha$ atau target significance level untuk melihat sub-lintasan mana yang tidak memiliki nilai $p$-value yang cukup 
\end{enumerate}

\section{Output}
Output dari Stat-DSM merupakan nilai significance level yang sudah dikoreksi $\delta*$.

\section{Test Case}
\subsection{Input}
\begin{table}[htbp]
    \centering
    \caption{$T$ Lintasan dengan $g$ Label}
    \label{tab:trajectory_data}
    \begin{tabular}{ccccc}
    \hline
    \textbf{TID} & \textbf{Label} & \textbf{Point 0} & \textbf{Point 1} & \textbf{Point 2} \\ \hline
    0   & 0     & (0, 0)      & (1, 0)      & (2, 0)      \\
    1   & 1     & (0.1, 0)    & (1.1, 0)    & (2.1, 0)    \\
    2   & 1     & (0, 0.1)    & (1, 0.1)    & (2, 0.1)    \\
    3   & 0     & (10, 12)    & (11, 12)    & (12, 12)    \\ \hline
\end{tabular}
\end{table}

\begin{enumerate}
    \item $\varepsilon = 0.5$
    \item $min_l = 2$
    \item $K = 1$
    \item $B = 4$
    \item $\alpha = 0.05$
\end{enumerate}

\subsection{Processing Awal}
\subsubsection{Buat permutasi label}
Karena $B = 4$, akan dibuat 4 permutasi dari label original.
Contoh :
\begin{table}[H]
    \centering
    \begin{tabular}{c|cccc}
        \toprule
        \textit{Permutasi ke-} & $T0$ & $T1$ & $T2$ & $T3$ \\
        \midrule
        1 & 0 & 0 & 1 & 1 \\
        2 & 0 & 1 & 0 & 1 \\
        3 & 1 & 0 & 0 & 1 \\
        4 & 1 & 0 & 1 & 0 \\
        \bottomrule
    \end{tabular}
\end{table}

\subsubsection{Hitung jumlah lintasan pada masing-masing kelas label}
Untuk label positif $N_+ = 2$, untuk label negatif $N_- = 2$.

\subsubsection{{Buat array list\_min\_p}}
Buat array list\_min\_p dengan ukuran $B$ untuk diisi dengan kumpulan $p$-value terkecil untuk setiap permutasi. 
Awal-awal isi seluruh array dengan $\alpha$ yaitu $0.5$. 

\subsection{Kalkulasi Stat-DSM}
\subsubsection{Generate kandidat pertama}
$S_{1,2}^{T0} = (0, 0) \rightarrow (1, 0)$

\subsubsection{Cari lintasan yang mirip}
Contoh bandingkan antara kandidat pertama
\[ A = (0, 0) \rightarrow (1, 0) \]
dengan sub-lintasan dari $T1$
\[ B =  (0.1, 0) \rightarrow (1.1, 0) \]
Menggunakan euclidean distance didapatkan :
\[ d_1=\sqrt{(0-0.1)^2+(0-0)^2} \]
\[ d_1=0.1 \]
dan
\[ d_2=\sqrt{(1-1.1)^2+(0-0)^2} \]
\[ d_2=0.1 \]
Dan karena $K = 1$ sehingga mengambil distance terbesar diantara keduanya 
\[ dist_1(A,B)=\max(0.1,0.1)=0.1 \]
Dan karena $\varepsilon = 0.5$ dan $0.1 \leq 0.5$ sehingga algoritma menanggap A mirip dengan B

Lakukan hal yang sama untuk semua sub-lintasan yang ditemukan seperti dari $T1$ masih terdapat $(1.1, 0) \rightarrow (2.1, 0)$ dan juga pada $T2$ terdapat $(0,0.1) \rightarrow (1, 0.1)$ dan $(1, 0.1) \rightarrow (2, 0.1)$ dan lain-lainnya untuk setiap lintasan.
Anggap didapatkannya :
\begin{table}[H]
    \centering
    \caption{Hasil Evaluasi Jarak Batas (Threshold)}
    \label{tab:jarak_threshold}
    \begin{tabular}{cccc}
        \hline
        \textbf{TID} & \textbf{Calculated distance} & \textbf{Threshold} & \textbf{Accepted?} \\ \hline
        $S_{1,2}^{T1}$   & 0.1                          & 0.5                & Yes                \\
        $S_{2,3}^{T1}$   & 2.1                          & 0.5                & No                 \\
        $S_{1,2}^{T2}$   & 0.1                          & 0.5                & Yes                \\
        $S_{2,3}^{T2}$   & 1.005                        & 0.5                & No                 \\
        $S_{1,2}^{T3}$   & Far greater than 0.5         & 0.5                & No                 \\
        $S_{2,3}^{T3}$   & Far greater than 0.5         & 0.5                & No                 \\ \hline
    \end{tabular}
\end{table}

\subsubsection{Kalkulasi Lower Bound}
Setelah mendapatkan semua sub-lintasan yang mirip dengan $S_1^{T0}$, kita perlu menghitung lower bound dari sub-lintasan tersebut untuk melihat apakah memperpanjang branch ini diperlukan.
Dimiliki $\text{support} = 4, N_+ = 2, N_- = 2$ sehingga didapat
Upper lower bound : 
\[ LB_U = \frac{\binom{2}{2}\binom{2}{0}}{\binom{4}{2}} \]
\[ LB_U = 0.1667 \]

dan Lower lower bound :
\[ LB_L = \frac{\binom{2}{0}\binom{2}{2}}{\binom{4}{2}} \]
\[ LB_L = 0.1667 \]

dan final lower bound didapat sebagai :
\[ LB = min(LB_U, LB_L) \]
\[ LB = 0.1667 \]

\subsubsection{Pruning}
Pada bagian ini akan digunakannya list\_min\_p = $[0.05, 0.05, 0.05, 0.05]$ dan parameter $B = 4$ dan $\alpha = 0.05$
untuk mendapatkan $idx = \lfloor B\alpha \rfloor$
\[ idx = \lfloor 4*0.05 \rfloor \]
\[ idx = 0 \]
Sehingga untuk pruning, akan digunakannya list\_min\_p yang di sort terlebih dahulu untuk mendapatkan sorted\_list\_min\_p.
Dan mengambil \[ sorted\_list\_min\_p[0] = 0.05 \]
Karena $0.05 \leq 0.1667$, maka untuk branch $S_1{T0}$ dapat di prune karena branch tidak dapat membuat $p$-value yang cukup rendah.

\section{Test Case 2}
\subsection{Input}
\subsubsection{Dataset}
\begin{table}[H]
    \centering
    \caption{Data Koordinat Berdasarkan TID dan Label}
    \label{tab:data_koordinat}
    \begin{tabular}{ccccc}
        \hline
        \textbf{TID} & \textbf{Label} & \textbf{Point 0} & \textbf{Point 1} & \textbf{Point 2} \\ \hline
        0 & 0 & (0, 0) & (1, 0) & (2, 0) \\ 
        1 & 1 & (0.1, 0) & (1.1, 0) & (2.1, 0) \\ 
        2 & 1 & (0, 0.1) & (1, 0.1) & (2, 0.1) \\ 
        3 & 1 & (0, -0.1) & (1, -0.1) & (2, -0.1) \\
        4 & 1 & (0.1, 0.1) & (1.1, 0.1) & (2.1, 0.1) \\
        5 & 0 & (10, 10) & (11, 10) & (12, 10) \\ 
        6 & 0 & (10, 11) & (11, 11) & (12, 11) \\ 
        7 & 0 & (10, 12) & (11, 12) & (12, 12) \\ \hline
    \end{tabular}
\end{table}

\subsubsection{Parameters}
\begin{enumerate}
    \item $\varepsilon = 0.5$
    \item $min_l = 2$
    \item $K = 1$
    \item $B = 1000$
    \item $\alpha = 0.05$
\end{enumerate}

\subsection{Initial Processing}
\[ N_+ = 4, N_- = 4 \]
\[ list\_min\_p = [0.05] * 1000 \]

Permutations :
\begin{table}[H]
    \centering
    \begin{tabular}{ccccccccc}
        \toprule
        \textit{Permutasi ke-} & $T0$ & $T1$ & $T2$ & $T3$ & $T4$ & $T5$ & $T6$ & $T7$ \\
        \midrule
        1 & 1 & 0 & 0 & 0 & 0 & 1 & 1 & 1 \\
        \vdots & \vdots & \vdots & \vdots & \vdots & \vdots & \vdots & \vdots & \vdots \\
        1000 & 0 & 1 & 0 & 0 & 1 & 1 & 1 & 0 \\ \hline
    \end{tabular}
\end{table}

\subsection{Kalkulasi Stat-DSM}
Anggap kandidat pertama merupakan $S_1^{T0} = (0,0) \rightarrow (1,0)$ \\
Anggap kita mendapatkan lintasan yang mirip merupakan lintasan \[[T1, T2, T3, T4]\]

\subsection{Hitung lower bound \& pruning}
Upper lower bound :
\[ LB_U = \frac{\binom{4}{4} \binom{4}{0}}{\binom{8}{4}} \]
\[ LB_U = \frac{1}{70}=0.014286 \]
Lower lower bound :
\[ LB_L = \frac{\binom{4}{0} \binom{4}{4}}{\binom{8}{4}} \]
\[ LB_L = \frac{1}{70}=0.014286 \]
Final lower bound :
\[ LB = 0.014286 \]
Dan bandingkan dengan $\alpha$
\[ 0.014286 < 0.05 \]
Lolos pruning

\subsection{Update list\_min\_p}
\subsubsection{Menghitung $p$-value kandidat dengan label original}
Upper $p$-value :
\[ P(X\geq4) = \frac{\binom{4}{4}\binom{4}{0}}{\binom{8}{4}} \]
\[ P(X\geq4) = \frac{1}{70} = 0.014286 \]
dan Lower $p$-value :
\[ P(X\leq4) = 1 \]
sehingga two sided $p$-value :
\[ p = \min(1,2\times\min(P_{upper},P_{lower})) \]
\[ p = 2\times0.014286 \]
\[ p = 0.028571 \]

\subsubsection{Permutation testing}
Pada bagian ini hanya akan dilakukan perhitungan untuk 5 permutasi, pada nyatanya dapat dilakukan lebih.
Seharusnya akan dilakukan \[ B=\binom{8}{4}=70 \text{ permutasi} \]
Sub-lintasan $S_1^{T0}$ memiliki kemiripan dengan $T = [1,2,3,4]$. \\
Anggap $X = \text{jumlah tetangga(lintasan dengan kemiripan dengan sub-lintasan) dengan label positif}$
Anggap didapati :
\begin{table}[h]
\centering
\caption{Permutasi Data Table}
\label{tab:permutasi_data1}
\begin{tabular}{lccccccccc}
\hline
\textbf{Permutasi} & \textbf{T0} & \textbf{T1} & \textbf{T2} & \textbf{T3} & \textbf{T4} & \textbf{T5} & \textbf{T6} & \textbf{T7} & \textbf{X} \\ \hline
P1        & 0  & 1  & 1  & 1  & 1  & 0  & 0  & 0  & 4 \\
P2        & 1  & 0  & 0  & 0  & 0  & 1  & 1  & 1  & 0 \\
P3        & 0  & 1  & 0  & 1  & 0  & 1  & 0  & 1  & 2 \\
P4        & 1  & 1  & 0  & 1  & 0  & 0  & 1  & 0  & 2 \\
P5        & 0  & 1  & 1  & 1  & 0  & 1  & 0  & 0  & 3 \\ \hline
\end{tabular}
\end{table}

Sehingga kita dapat menghitung $p$-value dengan menggunakan \[ P(X=x)=\frac{\binom{4}{x}\binom{4}{4-x}}{\binom{8}{4}} \]
Dan hasilnya :
\begin{table}[H]
\centering
\begin{tabular}{|c|c|c|}
\hline
\textbf{X} & \textbf{Perhitungan} & \textbf{Jumlah kemungkinan} \\ \hline
0 & $\binom{4}{0}\binom{4}{4}$ & 1 \\ 
1 & $\binom{4}{1}\binom{4}{3}$ & 16 \\
2 & $\binom{4}{2}\binom{4}{2}$ & 36 \\
3 & $\binom{4}{3}\binom{4}{1}$ & 16 \\
4 & $\binom{4}{4}\binom{4}{0}$ & 1 \\ \hline
\end{tabular}
\caption{Tabel Distribusi Kemungkinan X}
\label{tab:kemungkinan_x}
\end{table}

Selanjutnya harus dihitung upper-$p$-value dan lower-$p$-value
Didapatkan :
\begin{table}[h!]
\centering
\begin{tabular}{|c|c|c|}
\hline
\textbf{X} & \textbf{Upper-tail probability} & \textbf{Lower-tail probability} \\ \hline
0          & 1                               & 1/70                            \\ \hline
1          & 69/70                           & 17/70                           \\ \hline
2          & 53/70                           & 53/70                           \\ \hline
3          & 17/70                           & 69/70                           \\ \hline
4          & 1/70                            & 1                               \\ \hline
\end{tabular}
\caption{Tabel Probabilitas Ekor Atas dan Ekor Bawah}
\label{tab:probabilitas}
\end{table}

dan kemudian dapat dihitung $p$-value menggunakan rumus \[ p=\min(1,2\times\min(P_{upper},P_{lower})) \]
sehingga didapatkan nilai 
\begin{table}[h]
\centering
\caption{Tabel Nilai P-value dan Jumlah Permutasi}
\label{tab:permutasi_data2}
\begin{tabular}{ccc}
\hline
\textbf{X} & \textbf{Two-sided p-value} & \textbf{Jumlah permutasi} \\ \hline
0          & 0.028571                   & 1                         \\
1          & 0.485714                   & 16                        \\
2          & 1.000000                   & 36                        \\
3          & 0.485714                   & 16                        \\
4          & 0.028571                   & 1                         \\ \hline
\end{tabular}
\end{table}

\subsubsection{Ubah list\_min\_p}
Yang awalnya isi $list\_min\_p = [0.05] * 1000$
Kita akan mengubah semua nilai pada list\_min\_p yang memiliki angka yang lebih besar dibandingkan nilai Two-sided $p$-value yang kita bru hitung 
Sehingga sekarang $list\_min\_p = [0.028571,0.05,0.05,0.05,0.028571,\ldots,0.05]$

\subsection{Adjusted Significance Threshold}
Selanjutnya untuk mendapatkan significance threshold yang telah di adjust perlu dihitung
\[ \delta* = \lfloor B\alpha \rfloor\]
Angka $\delta*$ kemudian akan merujuk kepada sorted\_list\_min\_p sebagai index, dan value pada array tersebutlah menjadi adjusted significance threshold baru

\subsection{Loop}
Lakukan hal yang sama lagi dari atas dengan kandidat yang lain namun sorted\_list\_min\_p menggunakan yang baru
\end{document}